%! TEX root = ../am2-20-21.tex

\chapter{Recapitulare parțial}

\indent\indent \textbf{Integrale $ B, \Gamma $}

Folosind integralele Euler $ B, \Gamma $, calculați:

\begin{enumerate}[(a)]
    \item $ \ds\int_0^\infty \dfrac{x^2}{(1 + x^4)^2} dx $;
    \item $ \ds\int_0^\infty \dfrac{x}{x^6 + 1} dx $;
    \item $ \ds\int_0^\infty \dfrac{dx}{\sqrt[3]{x}(1 + x^2)} $;
    \item $ \ds\int_{-1}^\infty e^{-x^2 + 2x - 3} dx $.
\end{enumerate}

\vspace{0.3cm}

\textbf{Integrale duble}

1. Calculați integralele duble:

\begin{enumerate}[(a)]
    \item $ \ds\iint_D \sqrt{x^2 + y^2} dxdy $, unde $ D = \{ (x, y) \in \RR^2 \mid x^2 + y^2 \leq 16, x, y \geq 0 \} $;
    \item $ \ds\iint_D 3 - \sqrt{x^2 + y^2} dxdy $, unde $ D = \{ (x, y) \in \RR^2 \mid x^2 + y^2 \leq 9, x, y \geq 0 \} $.
\end{enumerate}

2. Calculați ariile domeniilor:
\begin{enumerate}[(a)]
    \item $ D = \{ (x, y) \in \RR^2 \mid x^2 \leq y \leq 3x^3, x \in [0,2] \} $;
    \item $ D = \{ (x, y) \in \RR^2 \mid 2x^3 \leq y \leq 5x^2, x \in [0, 1] \} $;
    \item $ D = \{ (x, y) \in \RR^2 \mid 0 \geq y \geq x^2 + 2x - 3, -1 \leq x \leq 1 \} $.
\end{enumerate}

\vspace{0.3cm}

\textbf{Integrale triple}

Calculați volumul mulțimilor $ \Omega \seq \RR^3 $, delimitate de cuadricele date:
\begin{enumerate}[(a)]
    \item $ z = 2x^2 + 2y^2 $ și $ \dfrac{9}{2} - z = x^2 + y^2 $;
    \item $ z = 3x^2 + 3y^2 $ și $ 2 - z = x^2 + y^2 $;
    \item $ 2z = x^2 + y^2 $ și $ 1 - z = 2x^2 + 2y^2 $.
\end{enumerate}

\vspace{0.3cm}

\textbf{Integrale curbilinii}

1. Aflați $ m \geq 0 $ astfel încît lungimea drumului parametrizat:
\[
    \gamma : \left[ \dfrac{\pi}{2}, \dfrac{3\pi}{2} \right] \to \RR^3, \quad %
    \gamma(t) = (t, m\cos t, -m\sin t)  
\]
să fie $ \pi \sqrt{10} $.

\vspace{0.3cm}

2. Fie funcția $ f : [0, 1] \to \RR, f(x) = 1 + x^2 $. Calculați masa
unui fir de densitate $ \rho(x) = x $, definit de graficul funcției date.

\vspace{0.3cm}

3. Calculați circulația cîmpului vectorial $ \vec{v} = (xy^2 + x^3, x + y^2) $
de-a lungul curbei:
\[
    C = \{ (x, y) \in \RR^2 \mid x^2 + y^2 = 4, x < 0 \} \cup %
    \{ (x, y) \in \RR^2 \mid x = 0, y \in [-2, 2] \}.
\]

\vspace{0.3cm}

4. Aceeași cerință ca la exercițiul 3 pentru cîmpul $ \vec{v} = (x^2y + y^3, xy + x) $,
de-a lungul curbei:
\[
    C = \{ (x, y) \in \RR^2 \mid x^2 + y^2 = 9, y < 0 \} \cup % 
    \{ (x, y) \in \RR^2 \mid y = 0, x \in [-3, 3] \}
\]